\documentclass[11pt]{article}
\usepackage[T1]{fontenc}
\usepackage[utf8]{inputenc}
\usepackage{lmodern}
\usepackage{amsmath,amssymb}
\usepackage{longtable,booktabs,array}
% Compatibility with Pandoc's captionless tables on older longtable versions.
\ifcsname c@none\endcsname\else\newcounter{none}\fi
\usepackage{graphicx}
\usepackage[margin=25mm]{geometry}
\usepackage{hyperref}
\hypersetup{colorlinks=true,linkcolor=blue,urlcolor=blue}
\providecommand{\tightlist}{\setlength{\itemsep}{0pt}\setlength{\parskip}{0pt}}
\providecommand{\pandocbounded}[1]{#1}
\setlength{\emergencystretch}{3em}
\setcounter{secnumdepth}{-1}
\begin{document}
\section{Complex analytic spaces and
analytification}\label{complex-analytic-spaces-and-analytification}

\emph{Written and self-checked by GPT-6.1 Sol (OpenAI) in Codex, Ultra
setting, October 2026. Analytic prerequisite integration by GPT-6 Astra
(OpenAI), Codex, Ultra. The linked analytic bridge was written and
self-checked by Claude Opus 5.5; a full independent review is not
claimed. The preparation and division proof in Section 2 is by Claude
Opus 5.5 (Anthropic). Original text: CC0.}

An equation with complex coefficients defines both a scheme and a space
of complex points. The second space has many more open sets and
functions: small balls are open, and convergent power series are
allowed. Nevertheless, at a complex point the two local rings have
exactly the same formal completion. This is the local reason that
passing from algebraic coherent sheaves to analytic coherent sheaves
preserves exact sequences. The next lesson will explain why projectivity
turns this local comparison into a global equivalence.

Throughout, \(X\) is a scheme locally of finite type over \(\mathbf C\).
It may have nilpotents, be nonseparated, and have infinitely many affine
charts. A \textbf{variety} will still mean a reduced separated scheme of
finite type. Analytification retains the entire structure sheaf,
including nilpotents. For nonseparated schemes we allow analytic locally
ringed spaces whose underlying space is non-Hausdorff; Hausdorffness and
global second countability are stated only with hypotheses that imply
them.

\subsection{1. Analytic equations and their
sheaves}\label{analytic-equations-and-their-sheaves}

For an open set \(W\subset\mathbf C^n\), let \(\mathcal O_W\) be the
sheaf of holomorphic functions. Iterated Cauchy integrals on polydiscs
give Taylor expansions; see
\href{../complex-analytic-spaces-and-coherent-sheaves/holomorphic-functions-of-several-variables.html\#1-polydiscs-and-the-definition-of-holomorphy}{Holomorphic
functions of several variables, Theorems 1.2 and 2.1} for the integral
and coefficient proofs. Demailly's freely available \emph{Complex
Analytic and Differential Geometry} is additional reading. At the origin
its stalk is

\[
H_n=\mathbf C\{z_1,\ldots,z_n\},
\]

the ring of power series converging in some neighbourhood of zero. Its
maximal ideal is \((z_1,\ldots,z_n)\). A germ with nonzero constant term
has a convergent reciprocal after shrinking the neighbourhood; every
other germ belongs to that ideal.

An \textbf{analytic subset} \(A\subset W\) is locally the common zero
set of finitely many holomorphic functions. Its vanishing ideal sheaf
\(\mathcal I_A\) consists of all germs vanishing on \(A\), and its
reduced structure sheaf is

\[
\mathcal O_A=(\mathcal O_W/\mathcal I_A)|_A.
\]

It embeds in the sheaf of continuous complex-valued functions on \(A\).
For the general construction choose a coherent ideal
\(\mathcal J\subset\mathcal O_W\), put \(A=V(\mathcal J)\), and use the
\textbf{quotient model}

\[
\bigl(A,(\mathcal O_W/\mathcal J)|_A\bigr).
\]

We retain \(\mathcal J\), rather than replacing it by \(\mathcal I_A\).
A complex analytic locally ringed space is a space locally isomorphic to
such models, with residue field \(\mathbf C\) at every point. A
Hausdorff such space is a complex space in the usual separated
convention. For example, \(\mathcal O_{\mathbf C}/(z^2)\) gives a
one-point model with ring \(\mathbf C[\epsilon]/(\epsilon^2)\), whereas
its reduced model has ring \(\mathbf C\).

An analytic morphism is locally given by holomorphic coordinate
representatives whose substitution carries the target defining ideal
into the source defining ideal. The representatives determine a pullback
on holomorphic germs and on their quotient sheaves. This includes the
images of nilpotent sections: a continuous point map alone does not
specify such a morphism. To check independence of representatives,
telescope one coordinate at a time in a convergent power series to
obtain

\[
F(u)-F(v)=\sum_j(u_j-v_j)G_j(u,v)
\]

with holomorphic \(G_j\) near \((c,c)\). Thus coordinate representatives
congruent modulo the source ideal give the same pullback. A local
homomorphism of these germ rings is determined by its coordinate images:
it preserves maximal-ideal powers, Taylor truncation determines the map
modulo every such power, and Krull intersection in the Noetherian target
ring makes the resulting equality exact. This also explains the
equivalence with morphisms of locally ringed spaces over \(\mathbf C\).

A module is \textbf{coherent} if it is locally finitely generated and
the relations among every finite family of sections are locally finitely
generated. Finite generation at one stalk is weaker than coherence on a
neighbourhood. That distinction is the content of Oka's theorem.

\textbf{Analytic foundation and reading order.} Preparation and division
are proved in Section 2. After that argument, read the following local
analytic results in \emph{Complex analytic spaces and coherent sheaves},
then return to Section 3. Their proofs concern analytic germs and
coherent sheaves, not GAGA, so none presupposes the comparison theorem
that this course will prove. Jean-Pierre Demailly's freely available
\href{https://www-fourier.univ-grenoble-alpes.fr/~demailly/manuscripts/agbook.pdf}{\emph{Complex
Analytic and Differential Geometry}}, together with the human sources
credited in the bridge, supplies further reading. The
\href{prerequisites.html}{prerequisite and proof guide} records the
exact scope and the remaining dependencies.

\begin{itemize}
\tightlist
\item
  \hyperref[local-analytic-algebra]{Preparation and division, Section
  2 below}: if a germ is regular of order \(r\) in the last variable, it
  is a unit times a monic degree-\(r\) polynomial in that variable whose
  other coefficients vanish at the origin. Division has a unique
  polynomial remainder of degree less than \(r\).
\item
  \href{../complex-analytic-spaces-and-coherent-sheaves/coherent-sheaves-and-okas-theorem.html\#2-oka-s-coherence-theorem}{Coherence
  of holomorphic functions, Theorem 2.1}: \(\mathcal O_W\) is coherent.
  Oka's theorem uses division to reduce analytic relations to finitely
  many polynomial coefficient relations in one fewer variable, and
  proves that the generators work near the point as well as at it.
\item
  \href{../complex-analytic-spaces-and-coherent-sheaves/analytic-germs-local-parametrization-and-the-nullstellensatz.html\#5-the-nullstellensatz-and-dimension}{The
  local analytic Nullstellensatz, Theorem 5.1}: for an ideal
  \(J\subset H_n\), the ideal of germs vanishing on its zero germ is
  \(\sqrt J\).
\item
  \href{../complex-analytic-spaces-and-coherent-sheaves/cartans-coherence-theorem-and-complex-spaces.html\#1-the-ideal-sheaf-of-an-analytic-set}{Coherence
  of analytic ideal sheaves, Theorem 1.1}: \(\mathcal I_A\) is coherent
  for every analytic subset. Cartan's theorem constructs uniform local
  equations using finite projections and coefficient interpolation; the
  statement includes singular analytic subsets.
\end{itemize}

The first coherence result is sufficient for quotient models by any
finitely generated analytic ideal; the Nullstellensatz and Cartan's
theorem additionally identify their reductions and vanishing ideals. The
bridge's
\href{../complex-analytic-spaces-and-coherent-sheaves/cartans-coherence-theorem-and-complex-spaces.html\#3-complex-spaces}{Cartan's
coherence theorem and complex spaces, Definition 3.1 and Theorem 3.2}
expressly includes arbitrary coherent ideals. Here is the module
argument, so passage to a nonreduced quotient introduces no new
coherence assumption.

\textbf{Proposition 1.1 (coherence on quotient models).} The structure
sheaf of every quotient model is coherent. Its coherent modules are
exactly its locally finitely presented modules. Kernels, images,
cokernels, extensions, tensor products and internal Hom of coherent
modules are coherent. If \(i:A\hookrightarrow W\) is the model
inclusion, a module on the model is coherent exactly when its
pushforward is coherent over \(\mathcal O_W\).

\textbf{Proof.} First work over a coherent ring sheaf \(\mathcal R\). A
finite type submodule of a coherent module is coherent, because its
relations are the same relations as in the larger module. For a morphism
\(\mathcal F\to\mathcal G\) of coherent modules, choose local generators
\(f_j\) of \(\mathcal F\). The finitely many local generators of the
relations among their images in \(\mathcal G\) give, by taking the
corresponding combinations of the \(f_j\), generators of the kernel. The
image is of finite type in \(\mathcal G\), and both kernel and image are
therefore coherent. To prove coherence of the cokernel, lift any finite
family of its sections locally to \(\mathcal G\) and adjoin generators
of the image. Relations among this enlarged family in \(\mathcal G\)
project onto the required relations in the cokernel, giving finite
generation. For an extension, lift generators of the quotient and adjoin
generators of the submodule. Relations first give relations in the
quotient; lift their finite generators, and then impose the finitely
generated relations among the resulting sections in the submodule. This
proves coherence of extensions, and in particular of finite direct sums
and finite free modules.

Consequently a finite presentation has coherent cokernel. Conversely,
generators of a coherent module and generators of their relations give a
finite presentation. Tensoring two finite presentations gives a finite
presentation of their tensor product. Applying internal Hom to a
presentation of \(\mathcal F\) realizes
\(\mathcal Hom(\mathcal F,\mathcal G)\) as a kernel between finite sums
of \(\mathcal G\), proving its coherence and the identification of its
stalk with the Hom of the stalks.

Now take \(\mathcal R=\mathcal O_W\) and
\(\mathcal B=\mathcal R/\mathcal J\). Oka's theorem and the kernel and
cokernel statements make \(\mathcal J\) and \(\mathcal B\) coherent over
\(\mathcal R\). A module coherent over \(\mathcal R\) and annihilated by
\(\mathcal J\) is coherent over \(\mathcal B\): the kernel of any map
\(\mathcal B^q\to\mathcal M\) is coherent over \(\mathcal R\), and its
finite generators also generate it over \(\mathcal B\). In particular
\(\mathcal B\) is a coherent ring sheaf. In the other direction a finite
presentation over \(\mathcal B\) is a cokernel between coherent
\(\mathcal R\)-modules, hence coherent over \(\mathcal R\). Finally
\(\mathcal B\) vanishes off the closed set \(A\). Restriction to \(A\)
and pushforward along \(i\) identify such sheaves and their stalks; off
\(A\) all stalks are zero. The preceding equivalence therefore gives the
assertion on the model. \(\square\)

\subsection{2. Local analytic algebra}\label{local-analytic-algebra}

\textbf{Preparation and division.} Write \(z=(z',w)\) with
\(z'\in\mathbf C^{n-1}\). Let \(f\in H_n\) satisfy \(f(0,w)=w^rv(w)\)
with \(v(0)\neq0\). Then \(f=uP\) for exactly one pair consisting of a
unit \(u\in H_n\) and a polynomial

\[
P=w^r+a_1(z')w^{r-1}+\cdots+a_r(z'),\qquad a_j\in H_{n-1},\quad a_j(0)=0 .
\]

Moreover, every \(g\in H_n\) can be written in exactly one way as
\(g=fq+R\), with \(q\in H_n\) and \(R\in H_{n-1}[w]\) of degree less
than \(r\) in \(w\).

\textbf{Proof.} Holomorphy in several variables means continuity
together with holomorphy in each variable separately, as in
\href{../complex-analytic-spaces-and-coherent-sheaves/holomorphic-functions-of-several-variables.html\#1-polydiscs-and-the-definition-of-holomorphy}{Holomorphic
functions of several variables, Definition 1.1}. The one-variable facts
used below are proved in
\href{../foundations-of-von-neumann-algebras/cauchy-s-theorem-for-cycles-and-its-consequences.html}{Cauchy's
theorem for cycles and its consequences}: interchange of integrals over
compact parameter intervals
(\href{../foundations-of-von-neumann-algebras/cauchy-s-theorem-for-cycles-and-its-consequences.html\#OA-FND-CT-07}{Lemma
0.1}), Goursat's theorem, Cauchy's theorem and Cauchy's formula on discs
(\href{../foundations-of-von-neumann-algebras/cauchy-s-theorem-for-cycles-and-its-consequences.html\#OA-FND-CT-02}{Theorems
2.1--2.3}), and Taylor expansion, Morera's theorem and the identity
theorem
(\href{../foundations-of-von-neumann-algebras/cauchy-s-theorem-for-cycles-and-its-consequences.html\#OA-FND-CT-03}{Theorems
3.2, 3.4 and 3.7}).

If \(r=0\), then \(f\) is a unit; take \(P=1\), \(u=f\), \(q=g/f\) and
\(R=0\). Let \(r\geq1\).

\emph{Integrals over a fixed circle.} Let \(\varphi(x,\zeta)\) be
continuous for \(x\) in an open subset of some \(\mathbf C^m\) and
\(|\zeta|=\rho\), and holomorphic in \(x\) for each \(\zeta\). Then
\(x\mapsto\int_{|\zeta|=\rho}\varphi(x,\zeta)\,d\zeta\) is holomorphic.
It is continuous by uniform continuity on compact sets. In each variable
separately, Lemma 0.1 exchanges a triangle integral with the circle
integral, Goursat's theorem makes the result zero, and Morera's theorem
gives holomorphy.

\emph{The zeros in each fibre.} Choose \(\rho>0\) and an open polydisc
\(D'\) around \(0\in\mathbf C^{n-1}\) such that \(f\) is holomorphic on
a neighbourhood of \(\overline{D'}\times\{|w|\leq\rho\}\),
\(f(0,w)\neq0\) for \(0<|w|\leq\rho\), and \(f(z',\zeta)\neq0\) for
\(z'\in D'\) and \(|\zeta|=\rho\). Once \(\rho\) satisfies the second
condition, the third is arranged by shrinking \(D'\), using continuity
and the compactness of the circle. Fix \(z'\in D'\). The zeros of
\(f(z',\cdot)\) inside the circle are isolated, hence finite in number.
Factoring them out by Taylor expansion gives
\(f(z',w)=\prod_i(w-c_i)^{m_i}k(w)\), where \(k\) has no zero on a disc
slightly larger than \(\{|w|\leq\rho\}\). Hence

\[
w^j\,\frac{\partial_wf(z',w)}{f(z',w)}=\sum_im_i\,\frac{w^j}{w-c_i}+w^j\,\frac{k'(w)}{k(w)} .
\]

Integrating over \(|w|=\rho\), Cauchy's theorem removes the last term
and Cauchy's formula evaluates the others:

\[
s_j(z')=\frac1{2\pi i}\int_{|\zeta|=\rho}\zeta^j\,\frac{\partial_wf(z',\zeta)}{f(z',\zeta)}\,d\zeta=\sum_im_ic_i^j\qquad(j\geq0).
\]

Each \(s_j\) is holomorphic on \(D'\) by the previous paragraph. The
integer \(s_0(z')\) is continuous, hence constant, and \(s_0(0)=r\). So
every fibre has exactly \(r\) zeros inside the circle, counted with
multiplicity. Call them \(b_1(z'),\ldots,b_r(z')\). Individual zeros
need not depend holomorphically on \(z'\): think of \(w^2-z_1\).

\emph{The polynomial.} For complex numbers \(b_1,\ldots,b_r\) put
\(S_j=\sum_lb_l^j\) and \(\prod_l(w-b_l)=\sum_{k=0}^rA_kw^{r-k}\),
\(A_0=1\). For \(|w|\) larger than every \(|b_l|\), geometric series
give \(\sum_l(w-b_l)^{-1}=\sum_{j\geq0}S_jw^{-j-1}\) with \(S_0=r\).
Multiplying by \(\prod_l(w-b_l)\) gives

\[
\sum_{k=0}^{r-1}(r-k)A_kw^{r-1-k}=\sum_{i=0}^r\sum_{j\geq0}A_iS_j\,w^{r-1-i-j}.
\]

After the substitution \(w=1/t\) and multiplication by \(t^{r-1}\), both
sides are power series in \(t\) near \(0\), so their coefficients agree.
The coefficient of \(t^k\) is \((r-k)A_k\) on the left, which is zero
for \(k=r\), and \(rA_k+S_1A_{k-1}+\cdots+S_kA_0\) on the right. Hence
\(kA_k=-(S_1A_{k-1}+\cdots+S_kA_0)\) for \(1\leq k\leq r\). Define
holomorphic functions on \(D'\) by \(a_0=1\) and
\(a_k=-(s_1a_{k-1}+\cdots+s_ka_0)/k\). The recursion just proved,
applied to the zeros \(b_l(z')\), whose power sums are the \(s_j(z')\),
shows that

\[
P(z',w)=w^r+a_1(z')w^{r-1}+\cdots+a_r(z')=\prod_{l=1}^r\bigl(w-b_l(z')\bigr),
\]

and \(P(0,w)=w^r\), so \(a_k(0)=0\).

\emph{The unit.} For fixed \(z'\), the functions \(f(z',\cdot)\) and
\(P(z',\cdot)\) vanish to the same order at each \(b_l(z')\). Hence
\(f/P\) extends holomorphically across these points, and the extension
is nowhere zero for \(|w|<\rho\). By Cauchy's formula the extension
equals

\[
u(z',w)=\frac1{2\pi i}\int_{|\zeta|=\rho}\frac{f(z',\zeta)}{P(z',\zeta)(\zeta-w)}\,d\zeta\qquad(|w|<\rho),
\]

which is holomorphic in \((z',w)\) because \(P(z',\zeta)\neq0\) on the
circle. It has no zeros, so \(1/u\) is holomorphic too. The continuous
functions \(f\) and \(uP\) agree off the finitely many \(b_l(z')\) in
each fibre, so \(f=uP\).

\emph{Uniqueness of \(P\) and \(u\).} If \(\sum_k|c_k|\tau^{-k}<1\),
then for \(|w|\geq\tau\) the lower terms of
\(w^r+c_1w^{r-1}+\cdots+c_r\) have modulus less than \(|w|^r\), so all
roots lie in \(\{|w|<\tau\}\). Let \(f=u_1P_1\) be another such
factorization. Setting \(z'=0\) shows that \(P_1\) also has degree
\(r\). Choose a product neighbourhood \(D''\times\{|w|<\tau_1\}\), with
\(\tau_1\leq\rho\), on which \(f=u_1P_1\) holds with \(u_1\) zero-free.
Then choose \(\tau<\tau_1\) and shrink \(D''\) so that, by the bound
just proved, all roots of \(P(z',\cdot)\) and of \(P_1(z',\cdot)\) lie
in \(\{|w|<\tau\}\) for \(z'\in D''\). For each such \(z'\), both
polynomials are monic of degree \(r\) with exactly the zeros of
\(f(z',\cdot)\) in \(\{|w|<\tau_1\}\) as roots, with the same
multiplicities. Hence \(P_1=P\), and then \(u_1=u\).

\emph{Division.} Shrink \(D'\) and \(\rho\) so that \(g\) is holomorphic
on a neighbourhood of \(\overline{D'}\times\{|w|\leq\rho\}\) and all
roots of \(P(z',\cdot)\) lie in \(\{|w|<\rho\}\). The divided difference
\(K(z';\zeta,w)=(P(z',\zeta)-P(z',w))/(\zeta-w)\) is a polynomial in
\(w\) of degree less than \(r\). Its coefficients are polynomials in
\(\zeta\) and the \(a_k(z')\). For \(|w|<\rho\) put

\[
Q(z',w)=\frac1{2\pi i}\int_{|\zeta|=\rho}\frac{g(z',\zeta)}{P(z',\zeta)(\zeta-w)}\,d\zeta,\qquad
R(z',w)=\frac1{2\pi i}\int_{|\zeta|=\rho}\frac{g(z',\zeta)\,K(z';\zeta,w)}{P(z',\zeta)}\,d\zeta .
\]

Both are holomorphic, and \(R\) is a polynomial in \(w\) of degree less
than \(r\) with holomorphic coefficients. Insert
\(1/(\zeta-w)=P(z',w)/(P(z',\zeta)(\zeta-w))+K(z';\zeta,w)/P(z',\zeta)\)
into Cauchy's formula for \(g(z',\cdot)\). This gives \(g=PQ+R\), so
\(g=fq+R\) with \(q=Q/u\).

If also \(g=fq'+R'\), then \(R-R'=f(q'-q)=Pu(q'-q)\) on a product
neighbourhood of the form above. For each \(z'\), the polynomial
\((R-R')(z',\cdot)\) of degree less than \(r\) vanishes at the \(r\)
roots of \(P(z',\cdot)\), counted with multiplicity, so it is zero. Then
\(u(q'-q)\) vanishes off finitely many points in each fibre, hence
everywhere, and \(q'=q\). \(\square\)

Bounds for division that hold uniformly for bounded dividends, with
explicit constants, are proved in
\href{../analytic-finiteness-and-preparation/weierstrass-preparation-and-division.html\#division-finite-generation-and-elementary-algebra}{Weierstrass
preparation and division}, Theorem 2. They are not needed in this
course.

\textbf{Lemma 2.1.} The ring \(H_n\) is Noetherian, and its maximal-adic
completion is \(\mathbf C[[z_1,\ldots,z_n]]\).

\textbf{Proof.} Induct on \(n\), beginning with \(H_0=\mathbf C\). For a
nonzero ideal \(J\), choose a nonzero \(f\in J\). A linear coordinate
change makes its first nonzero homogeneous part nonzero on the last
coordinate axis. Preparation then supplies a monic Weierstrass
polynomial \(P\in J\). Division identifies \(H_n/(P)\), as an
\(H_{n-1}\)-module, with the free module of polynomial remainders of
degrees \(0,\ldots,r-1\). By induction it is Noetherian. The image of
\(J\) is finitely generated in this quotient, so lifts of these
generators together with \(P\) generate \(J\). The zero ideal causes no
difficulty.

Modulo \((z)^m\), Taylor truncation identifies \(H_n/(z)^m\) with
polynomials of total degree less than \(m\). Taking the inverse limit
gives the full formal power-series ring. A convergent germ with every
Taylor coefficient zero is zero, so the Taylor map is injective.
\(\square\)

For any ideal \(J\subset H_n\), exact finite-module completion gives

\[
\widehat{H_n/J}=\mathbf C[[z]]/J\mathbf C[[z]].
\]

The precise earlier proof is
\href{../AG-CA/AG-CA-19.html\#3-noetherian-exactness-tensoring-and-flatness}{Completion,
Theorem 3.1}: over a Noetherian ring, Artin--Rees compares the induced
and intrinsic filtrations on a finite submodule, compatible lifts make
inverse limits exact, and a finite presentation identifies completed
modules with tensor products. Its Theorem 3.2 then proves flatness of
completion by the ideal criterion, with faithful flatness for
maximal-ideal completion of a local ring. These statements apply to
\(H_n\), which Lemma 2.1 has just proved Noetherian; they require no
reducedness.

\textbf{Lemma 2.2 (holomorphic inverse and implicit functions).} A
holomorphic map between open subsets of \(\mathbf C^n\) whose derivative
is invertible at a point has a holomorphic inverse on smaller
neighbourhoods. If \(r\) equations have Jacobian rank \(r\), their zero
set is locally a complex manifold of dimension \(n-r\).

\textbf{Proof.} Translate the point and make a linear change so the map
is \(F(z)=z+g(z)\), where \(g(0)=0\) and \(dg(0)=0\). On a sufficiently
small closed ball, \(\|dg\|\leq c<1\). For \(w\) in a smaller ball, the
map \(z\mapsto w-g(z)\) maps the first ball into itself and is a
contraction. Its successive iterates converge uniformly, with a
geometric bound, to the unique solution \(z(w)\). The iterates are
holomorphic in \(w\), and local uniform convergence preserves
holomorphicity by Cauchy's formula. The estimate also gives injectivity
of \(F\) on the first ball. For the implicit statement, add
complementary coordinates to the given equations so that the combined
map has invertible derivative, and apply the inverse statement.
\(\square\)

Local analytic dimension is the Krull dimension of the local analytic
algebra.
\href{../complex-analytic-spaces-and-coherent-sheaves/analytic-germs-local-parametrization-and-the-nullstellensatz.html\#4-the-local-parametrization-theorem}{Local
finite parametrization, Theorem 4.1}, and
\href{../complex-analytic-spaces-and-coherent-sheaves/analytic-germs-local-parametrization-and-the-nullstellensatz.html\#5-the-nullstellensatz-and-dimension}{Dimension
of analytic germs, Theorem 5.4}, give the finite-projection proof: after
adapted linear coordinates, an irreducible germ is finite over a
polydisc of the dimension of its local ring and is a finite covering
away from a discriminant. This also treats dimension component by
component.

\subsection{3. Constructing analytification with the full
ideal}\label{constructing-analytification-with-the-full-ideal}

For a finite type affine scheme
\(X=\operatorname{Spec}(\mathbf C[z]/J)\), give
\(X(\mathbf C)=V(J)\subset\mathbf C^n\) its ordinary topology and the
quotient sheaf

\[
\mathcal O_{X^{\mathrm{an}}}
=\bigl(\mathcal O_{\mathbf C^n}/J\mathcal O_{\mathbf C^n}\bigr)|_{V(J)}.
\]

The polynomial ideal \(J\) is finitely generated, so its image ideal is
coherent by Oka and Proposition 1.1. This is a quotient model for every
\(J\), including nonradical ideals. Its stalk at \(x\) is
\(\mathbf C\{z-x\}/J\mathbf C\{z-x\}\). The zero set specifies the
points and topology; the ideal specifies the entire local ring. For
reduced schemes the following additional argument recovers the reduced
construction.

\textbf{Lemma 3.1.} For a radical algebraic ideal \(J\) and a complex
point \(x\), its extension to the ambient convergent local ring is
radical and equals the analytic vanishing ideal of its zero germ.

\textbf{Proof.} Translate \(x\) to zero and put

\[
R=\mathbf C[z]_{(z)},\qquad H=\mathbf C\{z\},\qquad S=\mathbf C[[z]].
\]

Both ambient completions are \(S\). We prove directly that the
completion of every reduced finite type complex algebra at a complex
point is reduced. This avoids requiring a theorem about arbitrary Nagata
rings.

\emph{A domain over a polynomial ring.} Let \(E\) be a finite type
complex domain. The earlier programme proof
\href{../AG-CA/AG-CA-09.html\#3-normalization-that-respects-ideals}{Krull
dimension and Noether normalization, Corollary 3.2}, gives a finite
inclusion \(D=\mathbf C[u_1,\ldots,u_d]\subset E\). Write
\(K=\operatorname{Frac}D\) and \(L=\operatorname{Frac}E\). Localizing
the finite \(D\)-module \(E\) at \(D\setminus\{0\}\) gives a
finite-dimensional domain over \(K\), hence a field, namely \(L\). Thus
\(L/K\) is finite.

In characteristic zero every irreducible polynomial is separable: its
derivative is nonzero and has smaller degree, so the greatest common
divisor with the polynomial is one. We also record the primitive-element
step. For two separable generators \(\alpha,\beta\), list the roots
\(\alpha_i,\beta_j\) of their minimal polynomials in an algebraic
closure. Choose \(c\in\mathbf C\) avoiding the finitely many values for
which distinct pairs have equal sums \(\alpha_i+c\beta_j\). For
\(\gamma=\alpha+c\beta\), the polynomials \(f(\gamma-cT)\) and \(g(T)\),
over \(K(\gamma)\), have exactly the one common root \(\beta\). Their
monic greatest common divisor is \(T-\beta\), so the Euclidean algorithm
shows \(\beta\in K(\gamma)\), and then \(\alpha\in K(\gamma)\).
Repeating with the finite algebra generators of \(E\) produces
\(\theta\in E\) with \(L=K(\theta)\).

Let \(P\in K[T]\) be its monic minimal polynomial. Every algebra
generator of \(E\) is a polynomial in \(\theta\) over \(K\). Clearing
their finitely many denominators, the coefficients of \(P\), and a
Bezout identity for \(P,P'\), gives a nonzero \(h\in D\) such that

\[
E_h=D_h[T]/(P),\qquad
U P+V P'=1\quad\text{in }D_h[T].
\]

The quotient equality follows from monic division: a remainder of degree
less than \(\deg P\) cannot vanish at \(\theta\) unless all its
coefficients are zero. The Bezout identity is possible because \(P\) is
separable.

\emph{Complete the polynomial base.} Let \(x\) be a complex point of
\(\operatorname{Spec}E\), and let its contraction in \(D\) be
\(q=(u_1-a_1,\ldots,u_d-a_d)\). Put \(D_q=T_0\) and

\[
C_0=\widehat{T_0}=\mathbf C[[u-a]].
\]

The ring \(C_0\) is a domain: the lowest nonzero homogeneous parts of
two nonzero series have nonzero product. As a finite torsion-free
\(D\)-module, \(E\) embeds in \(D^r\) for some \(r\): choose a
\(K\)-basis of \(L\), write a finite generating list of \(E\) in that
basis, and multiply all coordinates by one common nonzero denominator.
Localize this injection at \(q\) and tensor with the flat completion
\(C_0\). The resulting finite algebra

\[
B_0=E\otimes_D C_0
\]

embeds as a \(C_0\)-module in \(C_0^r\), so is torsion-free. Its
localization at \(h\) is \(C_0[1/h][T]/(P)\). This monic quotient is
free over the domain \(C_0[1/h]\), so it embeds into its scalar
extension to \(F=\operatorname{Frac}C_0\). The Bezout identity makes
\(P\) squarefree over \(F\); factoring it into distinct irreducibles and
using Chinese remainders identifies \(F[T]/(P)\) with a finite product
of fields. Hence \((B_0)_h\) is reduced. A nilpotent of \(B_0\) must
therefore be annihilated by some power of \(h\). The polynomial \(h\) is
nonzero in \(C_0\), and torsion-freeness makes that nilpotent zero. Thus
\(B_0\) is reduced.

\emph{Separate the completed points.} The finite \(T_0\)-algebra \(E_q\)
has finitely many maximal ideals \(m_1,\ldots,m_s\), all above \(q\).
Here is the integral field test involved in the contraction assertion.
If a domain \(A_0\) has an integral extension which is a field, the
inverse of any nonzero \(a\in A_0\) satisfies a monic equation; multiply
that equation by \(a^{n-1}\) to express \(a^{-1}\) in \(A_0\). Hence
\(A_0\) is a field. Apply this to the inclusion of \(T_0/(m_i\cap T_0)\)
in \(E_q/m_i\). Since \(T_0\) is local, the contraction must be \(q\).
The maximal ideals therefore correspond to those of the
finite-dimensional complex fibre algebra. That algebra is Artinian,
since every strict descending chain of ideals lowers its vector-space
dimension. The exact Artinian radical and decomposition argument is
\href{../AG-CA/AG-CA-03.html\#4-what-an-artinian-ring-looks-like}{Noetherian
and Artinian rings, Lemma 4.1 and Theorem 4.2}. If \(N=\bigcap_i m_i\),
nilpotence of the radical of the fibre gives

\[
N^e\subset qE_q\subset N
\]

for some \(e\). These filtrations are cofinal. The ideals \(m_i^k\) are
pairwise comaximal: if \(a+b=1\) with \(a\in m_i,b\in m_j\), expansion
of \((a+b)^{2k-1}\) puts one in \(m_i^k+m_j^k\). Chinese remainders and
\(N^k=\bigcap_i m_i^k\) consequently give

\[
B_0=\widehat{E_q}_{\,qE_q}
\cong\varprojlim_k E_q/N^k
\cong\prod_i\varprojlim_k E_q/m_i^k
\cong\prod_i\widehat{E_{m_i}}.
\]

The first equality is finite-module completion, and the last holds
because every element outside \(m_i\) is already a unit modulo
\(m_i^k\). A factor of a reduced product ring is reduced. In particular
\(\widehat{E_x}\) is reduced.

For a general reduced finite type algebra \(E\), its finitely many
minimal primes have zero intersection, giving an injection
\(E\hookrightarrow\prod_i E/p_i\). Finiteness and the nilradical
assertion are proved in
\href{../AG-CA/AG-CA-03.html\#2-polynomial-rings-and-finite-geometric-descriptions}{Noetherian
and Artinian rings, Proposition 2.2}. Localize at the chosen complex
point and use exact completion on these finite modules. The completed
ring embeds in a finite product of reduced completed domain rings, with
zero factors omitted. It is therefore reduced as well.

Apply this conclusion to \(\mathbf C[z]/J\). Exact completion identifies
\(\widehat{R/JR}\) with \(S/JS\), so that ring is reduced. The
completion of \(H/JH\) is the same quotient. Faithful flatness of local
completion embeds the Noetherian local ring \(H/JH\) into it, making
\(H/JH\) reduced. The earlier analytic Nullstellensatz now identifies
the vanishing ideal with \(\sqrt{JH}=JH\). \(\square\)

\textbf{Theorem 3.2.} Every scheme locally of finite type over
\(\mathbf C\) has an analytification, uniquely characterized by these
finite type affine quotient charts. Scheme morphisms induce analytic
morphisms, compatibly with identities and composition. Analytification
commutes with products over \(\mathbf C\), open immersions and closed
immersions, including their full defining ideals. It therefore commutes
with locally closed immersions. A separated scheme has Hausdorff
analytification; a finite type scheme has second countable
analytification.

\textbf{Proof.} We check presentations, localizations, maps and gluing
separately, without making an ideal radical.

\emph{Adding coordinates.} Suppose \(A=\mathbf C[z]/J\) and
\(g_1,\ldots,g_m\in\mathbf C[z]\). Adding the coordinate classes \(g_j\)
gives the presentation

\[
A=\mathbf C[z,t]/(J,t_1-g_1(z),\ldots,t_m-g_m(z)).
\]

Its point set is the graph of \(g\), with a homeomorphism to \(V(J)\).
At \((x,g(x))\), substitution \(t=g(z)\) gives

\[
\frac{\mathbf C\{z-x,t-g(x)\}}{(J,t-g(z))}
\ \cong\ \frac{\mathbf C\{z-x\}}{J}.
\]

Indeed, for every holomorphic germ \(F\), the divided-difference
identity of Section 1 writes \(F(z,t)-F(z,g(z))\) in the ideal
\((t-g(z))\). This proves both the kernel assertion and the inverse
substitution. The germ maps come from sheaf maps on the graph, so the
sheaves are isomorphic. Given two finite polynomial presentations of the
same algebra, each coordinate class is a polynomial in the other list.
The combined list is obtained by adding coordinates to either list.
These graph isomorphisms thus give a canonical isomorphism between the
two analytic quotient charts. The isomorphism fixes every element of
\(A\); substitution shows that using three presentations gives the
cocycle identity.

\emph{Principal opens.} For \(h\in A\), represented by a polynomial, a
presentation of \(A_h\) is \(\mathbf C[z,v]/(J,vh-1)\). Its points map
homeomorphically to \(\{h\ne0\}\subset V(J)\), an ordinary open subset.
There \(h\) is an analytic unit. Substitution \(v=1/h\), with the same
divided-difference argument, identifies the quotient sheaf with the
restriction of the original chart sheaf. Changing the representative of
\(h\) modulo \(J\) has no effect on these sheaf maps. Since principal
opens form a basis of the Zariski topology of an affine scheme, this
verifies compatibility with arbitrary open overlaps.

\emph{Maps.} A map
\(\operatorname{Spec}A\to\operatorname{Spec}(\mathbf C[y]/K)\) sends
each \(y_j\) to a class represented by a polynomial \(g_j(z)\). The
relations give \(K(g(z))\subset J\). Holomorphic substitution therefore
gives a map between the analytic quotient charts. Two choices of
polynomial representatives differ by elements of \(J\), so the
divided-difference identity proves equality of their pullbacks on every
holomorphic germ, including nilpotents in the quotient. Fractional
coordinates are handled on principal opens by the preceding calculation.
Substitution is associative, so this construction respects composition
and identities.

\emph{Gluing.} Cover \(X\) by affine opens of finite type over
\(\mathbf C\). On an overlap, principal opens inside these charts and
the just-constructed regular transition maps give inverse analytic
isomorphisms of the quotient sheaves. They agree on smaller principal
opens, and the substitution cocycle is exactly the algebraic cocycle.
The ordinary chart topologies and their sheaves consequently glue. This
argument allows infinitely many charts and requires no separation
assumption. Each stalk is the local analytic quotient already specified.
Applying the map construction on source and target charts glues every
scheme morphism. A different affine cover has a common refinement; the
same canonical transition maps give a unique isomorphism compatible with
all quotient charts. This is the asserted uniqueness.

\emph{Products.} If \(A=\mathbf C[z]/J\) and \(B=\mathbf C[t]/K\), then

\[
A\otimes_{\mathbf C}B=\mathbf C[z,t]/(J(z),K(t)).
\]

Its points have the product topology on \(V(J)\times V(K)\). The
analytic product has precisely the quotient model by the sum of these
two ideals in the disjoint coordinate variables: given two analytic maps
from a quotient model, combine their coordinate representatives; the two
ideal-containment conditions are exactly the condition for this combined
map. The combined coordinate list is unique, by the germ substitution
rule, establishing the product universal property and the canonical
isomorphism. Gluing product charts proves the assertion for all locally
finite type schemes. Notice that this uses convergent series in both
variable lists, rather than an uncompleted tensor product of analytic
stalk rings.

\emph{Immersions and ideals.} An open immersion is restriction, by the
principal-open calculation. For a closed immersion with ideal
\(I\subset A\), choose polynomial representatives \(\widetilde I\) in
the ambient ring. Its chart is the quotient by \((J,\widetilde I)\), and
its analytic sheaf is the quotient of the original analytic sheaf by
\(I\mathcal O_{X^{\mathrm{an}}}\). Thus the ideal and all its powers are
retained. The support is closed because it is the common zero set of the
chosen generators. The constructions agree on overlaps, so the same
conclusion holds for every closed immersion and, by restriction, for
locally closed immersions.

Finally, on any product of affine chart domains the inverse image of the
diagonal of \(X(\mathbf C)\) is the point set of the algebraic diagonal.
If \(X\) is separated, that diagonal is a closed immersion, so the
analytic diagonal is closed in every such product by the preceding
calculation. Hence the glued space is Hausdorff. If \(X\) is of finite
type, finitely many affine charts suffice, and the union of their
countable bases is a countable basis. Neither conclusion is imposed for
a general \(X\). \(\square\)

\textbf{Corollary 3.3 (reduction).} There is a canonical identification

\[
(X_{\mathrm{red}})^{\mathrm{an}}=(X^{\mathrm{an}})_{\mathrm{red}}.
\]

For reduced \(X\), the quotient construction is the reduced analytic
construction, and every reduced locally closed subvariety has its usual
induced reduced analytic structure.

\textbf{Proof.} In a polynomial chart put \(J_0=\sqrt J\). Lemma 3.1
shows that \(J_0H\) is radical. It contains \(JH\), so
\(\sqrt{JH}\subset J_0H\). Conversely, each element of \(J_0\) has a
power in \(J\), so \(J_0H\subset\sqrt{JH}\). These inclusions prove
equality. The local analytic Nullstellensatz identifies that radical
with the vanishing ideal of the common zero set. Passing to quotients
and gluing proves the assertions. \(\square\)

There is a morphism of locally ringed spaces \(a:X^{\mathrm{an}}\to X\),
taking a complex point to its closed scheme point and taking regular
functions to holomorphic functions. Zariski opens pull back to ordinary
opens, because their complements are algebraic zero sets. The ordinary
topology is finer: it contains small balls that need not be Zariski
open.

\subsection{4. The comparison at a
point}\label{the-comparison-at-a-point}

\textbf{Theorem 4.1.} Let \(X\) be any scheme locally of finite type
over \(\mathbf C\), and let \(x\in X(\mathbf C)\). The local map

\[
A=\mathcal O_{X,x}\longrightarrow B=\mathcal O_{X^{\mathrm{an}},x}
\]

is faithfully flat, with residue field \(\mathbf C\), and induces an
isomorphism \(\widehat A\cong\widehat B\). It identifies every
maximal-ideal residue quotient. Consequently the dimensions and
regularity of the two local rings agree, including when \(X\) has
nilpotents.

\textbf{Proof.} Translate \(x\) to zero in a finite type affine chart
and put

\[
R=\mathbf C[z]_{(z)},\qquad H=\mathbf C\{z\},\qquad
S=\mathbf C[[z]].
\]

The quotient construction itself, without Lemma 3.1 and without
reduction, gives \(A=R/JR\) and \(B=H/JH\). Both are Noetherian. If
\(\mathfrak m\) and \(\mathfrak n\) are their maximal ideals, then, for
every \(k\geq1\), polynomial and Taylor truncation give the compatible
isomorphisms

\[
A/\mathfrak m^k
\ \cong\ \mathbf C[z]/(J+(z)^k)
\ \cong\ B/\mathfrak n^k.
\]

A polynomial denominator with nonzero constant term has a truncated
geometric-series inverse. Every convergent multiplier of a generator of
\(J\) can also be truncated modulo \((z)^k\). Thus both the ambient
rings and the image ideals agree in these quotients. Taking inverse
limits, or applying exact finite-module completion to the ambient
quotients, yields

\[
\widehat A=S/JS=\widehat B,
\]

and the map induced by \(A\to B\) is the identity on this common
quotient. In particular \(\mathfrak mB=\mathfrak n\) and the residue
field is \(\mathbf C\).

By
\href{../AG-CA/AG-CA-19.html\#3-noetherian-exactness-tensoring-and-flatness}{Completion,
Theorems 3.1--3.2}, maximal-ideal completion of a Noetherian local ring
is faithfully flat. For clarity, the logical route is exact completion
on finite modules, then injectivity of
\(I\otimes\widehat A\to\widehat A\) for every finite ideal, then the
flatness criterion proved in
\href{../AG-CA/AG-CA-07.html\#2-the-ideal-and-tor-criteria}{Tor and flat
modules, Theorem 2.1}. The nonzero residue quotient proves faithfulness;
the full criterion is
\href{../AG-CA/AG-CA-08.html\#1-detecting-nonzero-modules}{Faithful
flatness, Theorem 1.1}. All these inputs have precisely the Noetherian
or module hypotheses just stated.

Write \(C\) for the common completion. It is flat over \(A\) and
faithfully flat over \(B\). Given an injection of arbitrary
\(A\)-modules \(M'\to M\), let \(K\) be the kernel after tensoring with
\(B\). Flatness over \(B\) identifies \(C\otimes_BK\) with the kernel
after tensoring with \(C\); flatness over \(A\) makes that kernel zero.
Faithfulness over \(B\) makes \(K=0\). Hence \(B\) is flat over \(A\).
This descent uses arbitrary modules, so establishes flatness, not only
exactness on finite modules.

It is faithfully flat because it is a local flat map. Explicitly, for
\(0\ne m\in M\), its cyclic submodule \(Am=A/I\) has
\(I\subset\mathfrak m\), so \(B/IB\) surjects onto
\(B/\mathfrak n=\mathbf C\) and is nonzero. Flatness embeds this tensor
of the cyclic submodule into \(B\otimes_AM\). Thus tensoring with \(B\)
detects every nonzero module. Moreover the map \(M\to M\otimes_AB\)
becomes a split injection after tensoring once more with \(B\), the
retraction multiplying the two \(B\)-factors. Faithful flatness reflects
injectivity, so the original map is injective. Taking \(M=A\) and
\(M=A/I\) proves \(A\hookrightarrow B\) and \(IB\cap A=I\) for every
ideal.

A Noetherian local ring and its completion have the same dimension and
the same cotangent space over the residue field; the exact earlier proof
is
\href{../AG-CA/AG-CA-19.html\#4-local-dimension-and-regularity}{Completion,
Theorem 4.1}. It compares the lengths of the identical residue
quotients, hence the degrees of their Hilbert--Samuel polynomials, and
the first graded pieces of their maximal ideals. The common quotients
above therefore give equal dimensions and cotangent dimensions for \(A\)
and \(B\). Regularity is the equality of these two numbers, proving the
assertion. \(\square\)

A Zariski dense open \(U\subset X\) is also ordinarily dense in
\(X^{\mathrm{an}}\). First suppose \(X\) reduced. If \(Z=X\setminus U\),
with its reduced structure, contained an analytic neighbourhood of
\(x\), its relative analytic vanishing ideal at \(x\) would be zero.
Lemma 3.1 identifies that ideal with the extended algebraic ideal, and
faithful flatness contracts it to zero. The finite algebraic ideal then
vanishes on a Zariski neighbourhood of \(x\), so that neighbourhood lies
in \(Z\), contradicting density. For arbitrary \(X\), apply the reduced
argument to \(X_{\mathrm{red}}\); Corollary 3.3 identifies both
underlying analytic point spaces. This reduction is necessary: nilpotent
germs vanish at every point but need not be zero in the unreduced
structure sheaf.

\textbf{Proposition 4.2 (finite isomorphisms are detected
analytically).} Let \(f:Z\to X\) be a finite morphism between schemes
locally of finite type over \(\mathbf C\). If \(f^{\mathrm{an}}\) is an
analytic isomorphism, then \(f\) is an isomorphism.

\textbf{Proof.} Work over a finite type affine open of \(X\). For a
complex point \(x\), let \(A=\mathcal O_{X,x}\) and
\(D=(f_*\mathcal O_Z)_x\), a finite \(A\)-algebra. The fibre is a finite
\(\mathbf C\)-scheme; its closed points have residue field \(\mathbf C\)
and are exactly the analytic points above \(x\). The analytic
isomorphism gives exactly one such point \(z\). Every maximal ideal of
the integral algebra \(D\) lies above the maximal ideal of \(A\). Thus
\(D\) is local, with maximal ideal \(\mathfrak q\), and
\(D=\mathcal O_{Z,z}\). The finite-dimensional local fibre algebra
\(D/\mathfrak m_AD\) is Artinian, so its maximal ideal is nilpotent.
Hence \(\mathfrak q^e\subset\mathfrak m_AD\subset\mathfrak q\) for some
\(e\), and these two ideal-adic filtrations are cofinal.

Exact finite-module completion consequently identifies

\[
D\otimes_A\widehat A=\widehat D_{\mathfrak m_AD}
=\widehat D_{\mathfrak q}.
\]

Theorem 4.1 is functorial for local maps, since its identifications are
Taylor truncations of the original ring maps. The analytic stalk
isomorphism therefore says that
\(\widehat A\to\widehat D_{\mathfrak q}\) is an isomorphism. Faithful
flatness of \(A\to\widehat A\) detects that the finite \(A\)-module map
\(A\to D\) is an isomorphism. The kernel and cokernel of
\(\mathcal O_X\to f_*\mathcal O_Z\) are coherent and vanish at every
complex point. On a finite type affine chart any nonempty closed support
has a complex point, so both vanish. Finally a finite morphism is the
relative spectrum of its finite algebra \(f_*\mathcal O_Z\); this
algebra has just been identified with \(\mathcal O_X\). Thus \(f\) is an
isomorphism. \(\square\)

A merely bijective continuous map need not preserve local rings; the
cusp below shows why the ring hypothesis in this proposition matters.

\subsection{5. Coherent modules and
compactness}\label{coherent-modules-and-compactness}

Define

\[
\mathcal F^{\mathrm{an}}
=\mathcal O_{X^{\mathrm{an}}}\otimes_{a^{-1}\mathcal O_X}a^{-1}\mathcal F.
\]

\textbf{Proposition 5.1.} For every scheme \(X\) locally of finite type
over \(\mathbf C\), analytification is exact and faithful on coherent
modules, preserves coherence, and commutes with tensor products and
internal Hom of coherent modules. It commutes with pullback along scheme
morphisms and with pushforward along closed immersions. If
\(\mathcal I\subset\mathcal O_X\) is a coherent ideal, then
\(\mathcal I^{\mathrm{an}}\) identifies with the ideal it generates in
\(\mathcal O_{X^{\mathrm{an}}}\), and

\[
(\mathcal O_X/\mathcal I)^{\mathrm{an}}
=\mathcal O_{X^{\mathrm{an}}}/\mathcal I^{\mathrm{an}}.
\]

These assertions retain arbitrary nilpotents and ideal powers.

\textbf{Proof.} The analytic stalk is \(B\otimes_A\mathcal F_x\), where
\(A\to B\) is Theorem 4.1. Flatness proves exactness; it even proves
exactness for arbitrary module sheaves, without a coherence assumption.
A coherent algebraic sheaf has a finite presentation on each
sufficiently small affine chart because that chart is Noetherian. Its
analytification has the same finite matrix presentation over the
coherent analytic quotient sheaf. Proposition 1.1 makes its cokernel
coherent, including on nonreduced models. Tensor compatibility is
associativity of tensor products.

For internal Hom, use a presentation \(A^r\to A^s\to M\to0\). Then
\(\operatorname{Hom}_A(M,N)\) is the kernel of \(N^s\to N^r\). Flatness
preserves that kernel, giving

\[
B\otimes_A\operatorname{Hom}_A(M,N)
\cong\operatorname{Hom}_B(B\otimes_A M,B\otimes_A N).
\]

The sheaf map is therefore an isomorphism on all analytic stalks; the
identification of stalk Hom with module Hom follows from that same
finite presentation, rather than from a general assertion about
arbitrary sheaves. If \(\mathcal F^{\mathrm{an}}=0\), faithful flatness
makes \(\mathcal F_x=0\) at every complex point. A nonzero coherent
module restricted to some finite type affine chart has nonempty closed
support there. Such a support has a maximal ideal, and the algebraic
Nullstellensatz gives a complex point. Hence \(\mathcal F=0\). If a
coherent-module morphism analytifies to zero, apply this argument to its
coherent image, using exactness. This proves faithfulness even when
\(X\) is not quasi-compact.

For \(f:X\to Y\), the definitions of \(f^*\) and analytification give

\[
(f^*\mathcal G)^{\mathrm{an}}
\cong(f^{\mathrm{an}})^*(\mathcal G^{\mathrm{an}}),
\]

by associativity along the commuting diagram of algebraic and analytic
local ring maps. These maps commute because they were all constructed by
the same coordinate substitutions. No flatness of \(f\) is needed for
this right-exact tensor identity.

For a closed immersion \(i:Y\hookrightarrow X\), write \(A\to A/I\) in a
chart and \(B\to B/IB\) in its analytification. At a complex point of
\(Y\) the two candidate pushforward stalks are canonically

\[
B\otimes_A M
\ \cong\ (B/IB)\otimes_{A/I}M,
\]

where \(M\) is the stalk on \(Y\). Outside the closed set both stalks
vanish. These natural stalk maps prove
\((i_*\mathcal F)^{\mathrm{an}}\cong i^{\mathrm{an}}_*(\mathcal F^{\mathrm{an}})\).
Finally, tensor
\(0\to\mathcal I\to\mathcal O_X\to\mathcal O_X/\mathcal I\to0\).
Exactness identifies its first image with the generated ideal and its
cokernel with the displayed quotient. On stalks the extended ideal of
\(I^k\) is \((IB)^k\), so powers are retained as well. \(\square\)

\textbf{Proposition 5.2.} If \(X\) is a projective scheme over
\(\mathbf C\), then \(X^{\mathrm{an}}\) is compact, whether or not \(X\)
is reduced.

\textbf{Proof.} The map from the unit sphere in \(\mathbf C^{n+1}\) to
complex projective space, sending a unit vector to its line, is
continuous and surjective. Thus \(\mathbf P^n(\mathbf C)\) is compact.
It is Hausdorff, for instance by identifying a line with its orthogonal
rank-one projection matrix. A closed immersion
\(X\hookrightarrow\mathbf P^n\) analytifies to a closed analytic
quotient model by Theorem 3.2. Its underlying point set is consequently
closed and compact. Nilpotents change its sheaf, while this compactness
proof concerns its point set. \(\square\)

\textbf{Proposition 5.3 (properness and compactness).} A separated
scheme \(X\) of finite type over \(\mathbf C\) is proper if and only if
\(X^{\mathrm{an}}\) is compact.

\textbf{Proof.} First take \(X\) reduced. The exact earlier programme
input is
\href{../AG-GS/prerequisites/AG-MO/projective-morphisms-and-chows-lemma.html}{\emph{Projective
morphisms and Chow's lemma}, Theorem 4.1}, in \emph{Morphisms of
schemes}. Its proof constructs a proper surjection \(\pi:V\to X\) with
an immersion \(V\hookrightarrow\mathbf P^n_{\mathbf C}\). Its hypotheses
hold here: \(\operatorname{Spec}\mathbf C\) is Noetherian, and \(X\) is
separated and of finite type. In this affine-base case the affine chart
embeddings simply come from finite sets of algebra generators. The proof
takes the schematic closure of a common dense affine-chart intersection
in a product of projective spaces; it glues the inverse images of the
affine charts and proves properness on that target cover. It uses
neither analytification nor proper coherent direct images. Replace \(V\)
by its reduction and let \(Y\) be its reduced closure in projective
space. Then \(V\) is open and dense in the projective reduced scheme
\(Y\), which need not be irreducible.

The graph \(V\to X\times Y\) is closed. To check this, \(V\) is proper
over \(X\), \(X\times Y\) is separated over \(X\), and any \(X\)-map
from the former to the latter is proper: factor it as its closed graph
followed by the base change of \(V\to X\). The graph map here is also a
locally closed immersion because \(V\hookrightarrow Y\) is an open
immersion. Its image is therefore closed, and is the graph with its
induced reduced structure.

If \(X\) is proper, composition makes \(V\) proper over \(\mathbf C\).
Its immersion into projective space is proper by the same graph
factorization, hence has closed image. Thus \(V=Y\). Proposition 5.2
makes \(V^{\mathrm{an}}\) compact. The map \(\pi^{\mathrm{an}}\) is
surjective: for each complex point of \(X\), its nonempty finite type
fibre has a closed point with residue field \(\mathbf C\). Its
continuous image \(X^{\mathrm{an}}\) is compact.

Conversely, assume \(X^{\mathrm{an}}\) compact. The analytic graph is a
closed subset of \(X^{\mathrm{an}}\times Y^{\mathrm{an}}\) by Theorem
3.2. This product is compact, so projection makes \(V^{\mathrm{an}}\)
compact and hence closed in the Hausdorff space \(Y^{\mathrm{an}}\).
Section 4 makes it dense there. Consequently
\(V(\mathbf C)=Y(\mathbf C)\). A nonempty closed subset of the finite
type scheme \(Y\) has a complex point, so the complement of \(V\) is
empty and \(V=Y\). A proper surjective image of a proper scheme is
universally closed: after any base change, the image in the base of a
closed subset of the target is the image of its closed inverse image in
the proper source. The assumed finite type and separation of \(X\) now
make it proper.

For a general \(X\), the closed immersion
\(X_{\mathrm{red}}\hookrightarrow X\) is surjective after every base
change: its locally nilpotent ideal stays locally nilpotent. The same
inverse-image argument shows that \(X\) is universally closed exactly
when \(X_{\mathrm{red}}\) is; the reverse implication also follows by
composition with that closed immersion. Since \(X\) is already separated
and of finite type, properness of the two schemes is equivalent.
Corollary 3.3 identifies their underlying analytic spaces. Applying the
reduced result proves the assertion. \(\square\)

\subsection{6. Points, nilpotents and
singularities}\label{points-nilpotents-and-singularities}

Affine space analytifies to \(\mathbf C^n\), with all holomorphic
functions on ordinary open subsets. The projective line analytifies to
the Riemann sphere, whose two standard coordinates are related by
\(w=1/z\).

The double point \(\operatorname{Spec}\mathbf C[\epsilon]/(\epsilon^2)\)
analytifies to a one-point space with ring
\(\mathbf C\{z\}/(z^2)=\mathbf C[\epsilon]/(\epsilon^2)\). The identity
and the map sending \(\epsilon\) to \(2\epsilon\) induce the same map of
point sets but different analytic ring maps. This is why retaining the
defining ideal is essential. Likewise, gluing two copies of
\(\mathbf A^1\) along \(D(z)\) produces two analytic copies of
\(\mathbf C\) glued along \(\mathbf C^*\): every neighbourhood of either
origin meets every neighbourhood of the other. Its analytification is
non-Hausdorff, as permitted by the nonseparated construction. An
uncountable disjoint union of complex points is locally of finite type
and has discrete uncountable analytification, so global second
countability would also exclude schemes covered by Theorem 3.2.

For the cuspidal cubic \(y^2=x^3\), the parametrization
\(t\mapsto(t^2,t^3)\) is a continuous bijection
\(\mathbf C\to X^{\mathrm{an}}\). Away from the origin its inverse is
\(y/x\), and at the origin continuity follows from \(|t|^2=|x|\). It is
a homeomorphism. It is not an analytic isomorphism: at the cusp the
local ring is

\[
\mathbf C\{x,y\}/(y^2-x^3),
\]

of dimension one and cotangent dimension two. The relation has no linear
term. This agrees with the algebraic singularity by Theorem 4.1.
Moreover, substituting \((t^2,t^3)\) into a convergent germ never
produces a linear term in \(t\), so the inverse parameter is not
holomorphic at the cusp.

Finally, \(\exp:\mathbf C\to\mathbf C^*\) is holomorphic. An algebraic
morphism \(\mathbf A^1\to\mathbf G_m\) would send the invertible
coordinate of \(\mathbf G_m\) to a unit in \(\mathbf C[z]\). All such
units are nonzero constants. Thus the exponential is not algebraic.
Faithfulness of analytification does not imply fullness on nonproper
varieties.

\subsection{7. Exercises with solutions}\label{exercises-with-solutions}

\textbf{Exercise 7.1 (easy).} Identify the analytification of affine
space and explain why algebraic regular functions do not exhaust its
global analytic functions.

\textbf{Solution.} The ideal is zero in the chart construction, so the
underlying space is \(\mathbf C^n\) and its sheaf is
\(\mathcal O_{\mathbf C^n}\). For \(n\geq1\), the entire function
\(e^{z_1}\) is not polynomial: its derivatives of every order in \(z_1\)
are nonzero, whereas sufficiently high derivatives of a polynomial
vanish.

\textbf{Exercise 7.2 (easy).} Prove compactness of
\(\mathbf P^n(\mathbf C)\), including \(n=0\).

\textbf{Solution.} Every complex line meets the unit sphere. The
quotient map from that compact sphere is continuous and surjective, so
its image is compact. When \(n=0\), the quotient is one point. The
projection-matrix description in Proposition 5.2 also verifies
Hausdorffness.

\textbf{Exercise 7.3 (medium).} Compute the completions at zero of
\(\mathbf C[z]_{(z)}\) and \(\mathbf C\{z\}\), and the map between them.

\textbf{Solution.} Modulo \((z)^m\), both rings consist of the same
polynomials of total degree less than \(m\). In the algebraic
localization a denominator with nonzero constant term has a finite
geometric-series inverse modulo \((z)^m\). In the analytic ring Taylor
truncation gives the same inverse. Their compatible quotients therefore
have inverse limit \(\mathbf C[[z]]\), and the induced completion map is
the identity on every coefficient.

\textbf{Exercise 7.4 (medium).} Explain both exactness and detection of
zero for coherent modules under analytification. Then compute the
analytification of \(X=\operatorname{Spec}\mathbf C[t]/(t^3)\), its
reduction and its maximal-ideal quotients. Show that the closed double
point \(Y\hookrightarrow X\) defined by \((t^2)\) retains its ideal
after analytification.

\textbf{Solution.} At a complex point tensor with the faithfully flat
local algebra \(B\). Flatness preserves an exact sequence, so every
analytic stalk sequence is exact. Faithfulness implies that a finite
algebraic stalk with zero tensor is zero. If a coherent algebraic module
were nonzero, its closed support would contain a complex point,
contradicting this detection. Applying this to the coherent image of a
morphism proves detection of zero morphisms as well.

The point set is \(\{0\}\), and Taylor division gives
\(\mathcal O_{X^{\mathrm{an}},0}=\mathbf C\{t\}/(t^3)=\mathbf C[t]/(t^3)\).
Its reduction is the one-point ring \(\mathbf C\). The quotients by the
first, second and third powers of the maximal ideal have bases \(1\),
then \(1,t\), then \(1,t,t^2\); all later quotients equal the last one.
The ideal of \(Y^{\mathrm{an}}\) is the nonzero one-dimensional ideal
generated by \(t^2\), and the quotient is \(\mathbf C[t]/(t^2)\).
Replacing that ideal by the set-theoretic vanishing ideal \((t)\) would
incorrectly replace the double point by its reduction.

\textbf{Exercise 7.5 (medium).} Prove that no nonconstant algebraic
morphism \(\mathbf A^1\to\mathbf G_m\) exists. Why does this not
prohibit nonconstant holomorphic maps?

\textbf{Solution.} A ring map \(\mathbf C[u,u^{-1}]\to\mathbf C[z]\)
must send \(u\) to a polynomial unit, which is constant by degree
additivity in a product equal to one. Analytically, a nowhere-zero
entire function can be nonconstant; \(e^z\) and its reciprocal are both
entire. The analytic function ring has additional units.

\textbf{Exercise 7.6 (harder).} A homeomorphism need not remove a
singularity. Verify this on the cuspidal parametrization and locate the
failed inverse function.

\textbf{Solution.} The inverse is continuous off the cusp by the
quotient \(y/x\), and at the cusp by \(|t|=|x|^{1/2}\). The derivative
of the parametrization at zero is zero, so Lemma 2.2 does not apply.
More decisively, the analytic cusp ring embeds into \(\mathbf C\{t\}\)
by substitution, but its image has no term of degree one. Hence it
cannot contain the inverse coordinate \(t\). Its cotangent dimension is
two, whereas that of \(\mathbf C\{t\}\) is one; the homeomorphism is not
an isomorphism of locally ringed spaces.

\subsection{8. Proof providers and further
reading}\label{proof-providers-and-further-reading}

The definitions and comparison follow the mathematical framework of
J.-P. Serre,
\href{https://www.numdam.org/item/AIF_1956__6__1_0/}{\emph{Géométrie
algébrique et géométrie analytique}}, Annales de l'Institut Fourier 6
(1956), 1--42, freely readable through Numdam. The article is credited
as the historical source; the proofs used here are the ones supplied in
this course and the specified openly reusable sources.

The primary freely reusable human treatment is Demailly's
\href{https://www-fourier.univ-grenoble-alpes.fr/~demailly/manuscripts/agbook.pdf}{\emph{Complex
Analytic and Differential Geometry}}, version of 21 June 2012. Its
second chapter, section 2.A, treats preparation and division by C. L.
Siegel's contour method. The bridge's Oka proof, Theorem 2.1 and Lemma
2.2, develops the polynomial-relation reduction from Demailly II,
Theorem 3.19 and Lemma 3.20. Its local-parametrization and
Nullstellensatz proof develops adapted coordinates, the primitive
element and the discriminant before proving Theorem 5.1 for every ideal
of the convergent ring; its Cartan proof develops the
parameter-dependent interpolation before treating the general quotient
models. These are the actual programme arguments used in Section 1, with
their exact theorem locations linked there.

The algebraic inputs have earlier programme proofs in \emph{Commutative
algebra for geometry}: \textbf{Completion}, Theorems 3.1--3.3 and 4.1;
\textbf{Tor and flat modules}, Theorem 2.1; \textbf{Faithful flatness},
Theorems 1.1 and 2.2; \textbf{Noetherian and Artinian rings},
Proposition 2.2, Lemma 4.1 and Theorem 4.2; and \textbf{Krull dimension
and Noether normalization}, Corollary 3.2. The complex-point assertions
use \href{../AG-CA/AG-CA-06.html}{\textbf{The Nullstellensatz and
Jacobson rings}, Theorems 1.3 and 2.1}, and closed support uses
\href{../AG-CA/AG-CA-02.html}{\textbf{Localization, local properties and
support}, Theorem 5.1 and Proposition 6.1}. The domain and
semilocal-completion arguments needed to pass from reduced algebraic
charts to reduced analytic charts are written out in Lemma 3.1. The
properness criterion uses the actual earlier proof of \textbf{Projective
morphisms and Chow's lemma}, Theorem 4.1, over the Noetherian base
\(\operatorname{Spec}\mathbf C\). Sections 3--5 establish the full
nonreduced construction and local comparison for locally finite type
schemes, and the finite-isomorphism detection used in the next lesson.

The
\href{https://kokunoyumeto.github.io/stacks-zh-hans-cn/en/gaga.html\#gaga-proposition-analytification}{AI
Integrated Stacks Project GAGA chapter} is additional comparison reading
for its reduced-variety treatment. Its local and global statements carry
that convention; the nonreduced statements proved here use the explicit
quotient and completion arguments above. Exact provider identities and
consumer hypotheses are recorded in the
\href{prerequisites.json}{prerequisite record}. Independently authored
teaching material is CC0, and linked programme components retain their
stated terms.
\end{document}
